% ===================================================================
\section{Example API Interaction Sequences}
\label{app:api-examples}
% ===================================================================

This appendix provides concrete examples of \API\ interactions covering the five
operation groups defined in Section~4.6.

\subsection{Example 1: \REO\ Submits a New \DPP}

An economic operator (EO) submits a new \DPP\ for a battery batch. The
interaction proceeds as follows:

\begin{enumerate}
\item \textbf{Authenticate}: The \EO\ presents its \eIDAS~2.0 credential to
obtain a bearer token.
\item \textbf{Submit DPP}: The \EO\ submits the signed \JSONLD\ \DPP\ via POST
\texttt{/dpps}.
\item \textbf{Ingestion}: The portal controller validates the \DPP\ against
\SHACL\ shapes, verifies the cryptographic signature, normalises external entity
references, and commits the \DPP\ to the triple-store.
\item \textbf{Register}: The \EO\ registers the \DPP\ with the EC central
registry via POST \texttt{/dpps/\{dppId\}/register}.
\end{enumerate}

\begin{lstlisting}[language=bash,breaklines=true,breakatwhitespace=true]
# Step 1: Authenticate
curl -X POST https://portal.example.eu/auth/token \
  -H "Content-Type: application/json" \
  -d '{"eidasCredential": "...",
       "requestedScopes": ["dpp:write"]}'

# Response: {"accessToken": "<JWT>", "attributes": {...}}

# Step 2: Submit DPP
curl -X POST https://portal.example.eu/dpps \
  -H "Authorization: Bearer <JWT>" \
  -H "Content-Type: application/ld+json" \
  -d @dpp-batch-2026-001.jsonld

# Response: {"dppId": "urn:dpp:product:batch-2026-001",
#            "status": "accepted",
#            "validationReport": {...}}

# Step 3: Register with EC central registry
curl -X POST \
  https://portal.example.eu/dpps/\
urn:dpp:product:batch-2026-001/register \
  -H "Authorization: Bearer <JWT>"

# Response: {"registryStatus": "registered",
#            "productIdentifier": \
#            "EN18215:DE:123456789:batch-2026-001"}
\end{lstlisting}

\subsection{Example 2: Consumer Scans \QR\ Code and Reads Public \DPP}

A consumer scans a \QR\ code on a product and requests the public-facing \DPP\
data. The interaction proceeds as follows:

\begin{enumerate}
\item \textbf{Authenticate}: The consumer's app presents a lightweight
\eIDAS~2.0 credential (no organisational identity required) to obtain a
\texttt{dpp:PublicUser} bearer token.
\item \textbf{Resolve \DPP}: The consumer's app requests the DPP via GET
\texttt{/dpps/\{dppId\}}.
\item \textbf{Filter}: The portal controller evaluates ODRL policies against the
\texttt{dpp:PublicUser} attributes and returns only the public-facing branches
(product name, basic composition summary, environmental impact indicators,
repairability information, recycling instructions).
\end{enumerate}

\begin{lstlisting}[language=bash,breaklines=true,breakatwhitespace=true]
# Step 1: Authenticate (lightweight)
curl -X POST https://portal.example.eu/auth/token \
  -H "Content-Type: application/json" \
  -d '{"eidasCredential": "...",
       "requestedScopes": ["dpp:read"]}'

# Response: {"accessToken": "<JWT>",
#            "attributes": {
#              "userCategory": "dpp:PublicUser", ...}}

# Step 2: Resolve DPP
curl -X GET \
  https://portal.example.eu/dpps/\
urn:dpp:product:batch-2026-001 \
  -H "Authorization: Bearer <JWT>" \
  -H "Accept: application/ld+json"

# Response: JSON-LD with only public-facing branches
# (dpp:filteredBranches: ["dpp:publicBranch"])
\end{lstlisting}

\subsection{Example 3: Circular Operator Traverses Supply Chain}

A certified recycler needs the full material composition and disassembly
information for a product at end-of-life. The interaction proceeds as follows:

\begin{enumerate}
\item \textbf{Authenticate}: The recycler presents its \eIDAS~2.0 credential
(including waste management permits and recycling certifications as \WthreeC\
Verifiable Credentials) to obtain a \texttt{dpp:CircularOperator} bearer token.
\item \textbf{Resolve \DPP}: The recycler requests the \DPP via GET
\texttt{/dpps/\{dppId\}}. The portal returns a filtered sub-graph including the
full composition data, hazard statements, and disassembly instructions.
\item \textbf{Traverse Supply Chain}: The recycler requests the full supply
chain via GET \texttt{/dpps/\{dppId\}/supplychain}. The portal recursively
follows \texttt{dpp:hasPrecursor} links and returns an aggregated JSON-LD graph
containing all accessible precursor \DPPs, each filtered according to the
recycler's access rights.
\end{enumerate}

\begin{lstlisting}[language=bash,breaklines=true,breakatwhitespace=true]
# Step 1: Authenticate
curl -X POST https://portal.example.eu/auth/token \
  -H "Content-Type: application/json" \
  -d '{
    "eidasCredential": "...",
    "verifiableCredentials": [
      "urn:vc:waste-mgmt-permit-DE-123",
      "urn:vc:recycling-cert-EU-456"
    ],
    "requestedScopes": ["dpp:read"]
  }'

# Response: {"accessToken": "<JWT>",
#            "attributes": {
#              "userCategory": "dpp:CircularOperator",
#              ...}}

# Step 2: Resolve DPP
curl -X GET \
  https://portal.example.eu/dpps/\
urn:dpp:product:batch-2026-001 \
  -H "Authorization: Bearer <JWT>" \
  -H "Accept: application/ld+json"

# Response: JSON-LD with composition, hazard,
#   and disassembly branches
# (dpp:filteredBranches: ["dpp:public",
#   "dpp:composition", "dpp:hazard",
#   "dpp:disassembly"])

# Step 3: Traverse full supply chain
curl -X GET \
  https://portal.example.eu/dpps/\
urn:dpp:product:batch-2026-001/supplychain \
  -H "Authorization: Bearer <JWT>" \
  -H "Accept: application/ld+json"

# Response: Aggregated JSON-LD graph with all
#   accessible precursor DPPs
\end{lstlisting}

\subsection{Example 4: Market Surveillance Authority Runs Cross-Product \SPARQL\
Query}

A market surveillance authority needs to identify all products in the portal
that contain a specific substance of concern. The interaction proceeds as
follows:

\begin{enumerate}
\item \textbf{Authenticate}: The authority presents its \eIDAS~2.0 governmental
credential to obtain a \texttt{dpp:MarketSurveillanceAuthority} bearer token.
\item \textbf{\SPARQL\ Query}: The authority submits a \SPARQL\ CONSTRUCT query
via POST \texttt{/sparql} that traverses all \DPPs\ in the portal and identifies
those containing the target substance. The portal controller executes the query
with full access rights (no \ODRL\ filtering for market surveillance
authorities) and returns the aggregated result.
\end{enumerate}

\begin{lstlisting}[language=bash,breaklines=true,breakatwhitespace=true]
# Step 1: Authenticate
curl -X POST https://portal.example.eu/auth/token \
  -H "Content-Type: application/json" \
  -d '{"eidasCredential": "<gov credential>",
       "requestedScopes": ["dpp:read","dpp:sparql"]}'

# Response: {"accessToken": "<JWT>",
#            "attributes": {
#              "userCategory": "dpp:MarketSurveillance",
#              ...}}

# Step 2: SPARQL query
curl -X POST https://portal.example.eu/sparql \
  -H "Authorization: Bearer <JWT>" \
  -H "Content-Type: application/json" \
  -d '{
    "query": "CONSTRUCT { ?dpp \
        dpp:containsSubstance ?substance . \
        ?dpp dpp:productName ?name } \
      WHERE { ?dpp \
        dpp:containsSubstance \
        <urn:echa:substance:7440-44-0> . \
        ?dpp dpp:productName ?name }",
    "format": "application/ld+json"
  }'

# Response: JSON-LD graph with all DPPs
#   containing the target substance
\end{lstlisting}

\subsection{Example 5: \REO\ Updates a \DPP\ After a Hazard Reclassification}

\ECHA\ reclassifies a substance, adding it to the \SVHC\ candidate list. The
\REO\ of a product containing that substance updates the \DPP\ to reflect the
new classification. The interaction proceeds as follows:

\begin{enumerate}
\item \textbf{Authenticate}: The \REO\ presents its \eIDAS~2.0 credential to
obtain a \texttt{dpp:InternalDataOriginator} bearer token.
\item \textbf{Retrieve \DPP}: The \REO\ retrieves the current \DPP\ via GET
\texttt{/dpps/\{dppId\}} (with full access rights as the data originator).
\item \textbf{Update \DPP}: The \REO\ submits the updated \DPP\ via PUT
\texttt{/dpps/\{dppId\}} with the corrected substance reference (now linking to
the updated \ECHA\ entry via the normalised \IRI). The portal controller
re-validates the \DPP\ against \SHACL\ shapes, verifies the signature, and
commits the updated version to the triple-store. All parent product \DPPs\ that
reference this \DPP\ via \texttt{dpp:hasPrecursor} automatically reflect the
updated classification upon subsequent resolution.
\end{enumerate}

\begin{lstlisting}[language=bash,breaklines=true,breakatwhitespace=true]
# Step 1: Authenticate
curl -X POST https://portal.example.eu/auth/token \
  -H "Content-Type: application/json" \
  -d '{"eidasCredential": "...",
       "requestedScopes": ["dpp:read","dpp:write"]}'

# Response: {"accessToken": "<JWT>",
#            "attributes": {
#              "userCategory": "dpp:InternalDataOriginator",
#              ...}}

# Step 2: Retrieve current DPP
curl -X GET \
  https://portal.example.eu/dpps/\
urn:dpp:product:batch-2026-001 \
  -H "Authorization: Bearer <JWT>" \
  -H "Accept: application/ld+json"

# Step 3: Update DPP with corrected substance ref
curl -X PUT \
  https://portal.example.eu/dpps/\
urn:dpp:product:batch-2026-001 \
  -H "Authorization: Bearer <JWT>" \
  -H "Content-Type: application/ld+json" \
  -d @dpp-batch-2026-001-updated.jsonld

# Response: {"dppId": "urn:dpp:product:batch-2026-001",
#            "status": "updated",
#            "validationReport": {...},
#            "version": 2}
\end{lstlisting}
